% ─────────────────────────────────────────────────────────────────────────────
%  Capítulo 2 – Arquitectura del compilador
% ─────────────────────────────────────────────────────────────────────────────
\chapter{Arquitectura del compilador}
\label{chap:arquitectura}

\section{Vista global del pipeline}

El compilador sigue la arquitectura clásica en fases consecutivas. Cada fase
transforma una representación del programa en otra de mayor nivel de abstracción
hacia la máquina destino, o bien extrae información (tabla de símbolos, tipos)
que se usa en fases posteriores.

\begin{figure}[H]
\centering
\begin{tikzpicture}[
  node distance=0.6cm and 1.8cm,
  box/.style={rectangle, draw, rounded corners=3pt,
              minimum width=3.8cm, minimum height=0.8cm,
              align=center, font=\small\ttfamily, fill=blue!8},
  arr/.style={-Stealth, thick},
  lbl/.style={font=\scriptsize\itshape, text=gray}
]

\node[box] (src)    {Fuente .java};
\node[box, below=of src]  (lex)    {Analizador Lexico\\(lexer.l)};
\node[box, below=of lex]  (par)    {Analizador Sintactico\\(parser.y)};
\node[box, below=of par]  (sem)    {Analizador Semantico\\(SemanticVisitor)};
\node[box, below=of sem]  (ir)     {Generacion de RI\\(IRBuilder)};
\node[box, below=of ir]   (opt)    {Optimizacion RI\\(IROptimizer)};
\node[box, below=of opt]  (reg)    {Asig. de Registros\\(RegisterAllocator)};
\node[box, below=of reg]  (x86)    {Generacion x86-64\\(X86Generator)};
\node[box, below=of x86]  (asm)    {Ensamblador .s};

\draw[arr] (src) -- (lex)
  node[lbl, right, midway] {codigo fuente};
\draw[arr] (lex) -- (par)
  node[lbl, right, midway] {tokens};
\draw[arr] (par) -- (sem)
  node[lbl, right, midway] {AST};
\draw[arr] (sem) -- (ir)
  node[lbl, right, midway] {AST tipado};
\draw[arr] (ir) -- (opt)
  node[lbl, right, midway] {cuadruplas IR};
\draw[arr] (opt) -- (reg)
  node[lbl, right, midway] {IR optimizado};
\draw[arr] (reg) -- (x86)
  node[lbl, right, midway] {RegMap};
\draw[arr] (x86) -- (asm)
  node[lbl, right, midway] {AT\&T asm};

% Rama de verificacion
\node[box, right=2.5cm of opt, fill=green!8, align=center]
  (wpgen)  {VCGenerator\\WPCalculus};
\node[box, below=0.6cm of wpgen, fill=green!8, align=center]
  (smt)    {SMTExporter\\(.smt2)};
\node[box, below=0.6cm of smt, fill=green!8]
  (z3)     {Z3 Solver};

\draw[arr, dashed, green!60!black]
  (sem) -| node[lbl, above, midway] {--verify} (wpgen);
\draw[arr, green!60!black] (wpgen) -- (smt);
\draw[arr, green!60!black] (smt) -- (z3);

\end{tikzpicture}
\caption{Pipeline completo del compilador JavaSubset. La rama verde
  (verificacion formal) se activa con la opcion \texttt{--verify}.}
\label{fig:pipeline}
\end{figure}

\section{Representaciones intermedias}

A lo largo del pipeline el programa adopta cuatro representaciones distintas:

\begin{enumerate}
  \item \textbf{Codigo fuente}: texto Java con posibles anotaciones formales.
  \item \textbf{AST} (\emph{Abstract Syntax Tree}): arbol de nodos tipados,
    jerarquia de clases C++ visitables mediante el patron \emph{Visitor}.
  \item \textbf{Cuadruplas IR}: secuencia plana de instrucciones de la forma
    $(op, \mathit{arg}_1, \mathit{arg}_2, \mathit{res})$, independiente de la
    arquitectura.
  \item \textbf{Ensamblador x86-64}: instrucciones AT\&T con registros reales
    o con posiciones en la pila.
\end{enumerate}

\section{Modos de operacion}

El compilador admite tres modos de operacion seleccionables mediante flags:

\begin{table}[H]
  \centering
  \caption{Flags del compilador y su efecto}
  \label{tab:flags}
  \begin{tabular}{llp{7cm}}
    \toprule
    \textbf{Flag}        & \textbf{Salida}   & \textbf{Descripcion} \\
    \midrule
    (ninguno)            & \texttt{prog.s}   & Compilacion estandar a ensamblador \\
    \texttt{--ir}        & \texttt{stderr}   & Imprime cuadruplas IR y forma lineal \\
    \texttt{--reg-alloc} & \texttt{stderr}   & Imprime mapa de asignacion de registros \\
    \texttt{--verify}    & \texttt{prog.smt2}& Verificacion formal con Z3 \\
    \bottomrule
  \end{tabular}
\end{table}

\section{Patron Visitor}

Todas las fases que procesan el AST utilizan el patron de diseno
\textbf{Visitor} (GoF). La interfaz abstracta \texttt{Visitor} declara un
metodo \texttt{visit()} para cada tipo de nodo del AST; cada fase
implementa esta interfaz sin modificar las clases del arbol.

\begin{figure}[H]
\centering
\begin{tikzpicture}[
  iface/.style={rectangle, draw, fill=yellow!15,
                minimum width=4cm, minimum height=0.7cm,
                align=center, font=\small},
  impl/.style={rectangle, draw, fill=blue!8,
               minimum width=3.5cm, minimum height=0.7cm,
               align=center, font=\small},
  arr/.style={-Stealth, thick}
]
\node[iface] (v) {\texttt{Visitor} (abstract)};
\node[impl, below left=1cm and 2.5cm of v] (sv) {\texttt{SemanticVisitor}};
\node[impl, below left=1cm and 0cm of v]   (ib) {\texttt{IRBuilder}};
\node[impl, below right=1cm and 0cm of v]  (av) {\texttt{ASTVisualizer}};
\node[impl, below right=1cm and 2.5cm of v] (dv) {\texttt{DOTVisualizer}};

\draw[arr] (sv) -- (v);
\draw[arr] (ib) -- (v);
\draw[arr] (av) -- (v);
\draw[arr] (dv) -- (v);
\end{tikzpicture}
\caption{Jerarquia del patron Visitor en el compilador.}
\label{fig:visitor}
\end{figure}

\section{Estructura de directorios}

\begin{lstlisting}[style=pseudo, caption={Organizacion del repositorio}]
.
|-- include/         Cabeceras: ASTNodes.h, IR.h, Predicate.h, ...
|-- src/             Implementaciones: lexer.l, parser.y, *Visitor.h, ...
|   |-- lexer.l           Analizador lexico (Flex)
|   |-- parser.y          Gramatica + acciones semanticas (Bison)
|   |-- SemanticVisitor.h Comprobacion de tipos
|   |-- IRBuilder.h       Traduccion AST -> IR
|   |-- IROptimizer.h     Optimizaciones sobre IR
|   |-- RegisterAllocator.h  Chaitin-Briggs coloring
|   |-- X86Generator.h    Generacion de codigo x86-64
|   |-- WPCalculus.h      Calculo de precondicion mas debil
|   |-- VCGenerator.h     Generacion de condiciones de verificacion
|   `-- SMTExporter.h     Exportacion a SMT-Lib2 y ejecucion de Z3
|-- tests/
|   |-- backend/          Programas de prueba end-to-end
|   |-- semantic/         Pruebas de deteccion de errores
|   |-- verification/     Demos de verificacion formal
|   `-- run_tests.sh      Script de pruebas
`-- Makefile
\end{lstlisting}
